\documentclass[french]{article}
\usepackage[utf8]{inputenc}
\usepackage[T1]{fontenc}
\usepackage{lmodern}
\usepackage{geometry}
\usepackage{graphicx}
\usepackage{babel}
\usepackage{amsmath}
\usepackage{amsthm}
\usepackage{amssymb}
\usepackage{hyperref}
\usepackage{stmaryrd}
\usepackage{enumitem}
\usepackage{tcolorbox}
\usepackage{dsfont}
\usepackage{multirow}
\usepackage{etoc}
\usepackage{titlesec}
\usepackage{esint}
\usepackage{cancel}
\usepackage{mathdots}

\setlist[itemize]{label=$\bullet$}


\renewcommand{\P}{\mathbb{P}}
\newcommand{\N}{\mathbb{N}}
\newcommand{\Z}{\mathbb{Z}}
\newcommand{\Q}{\mathbb{Q}}
\newcommand{\R}{\mathbb{R}}
\newcommand{\C}{\mathbb{C}}
\newcommand{\K}{\mathbb{K}}
\renewcommand{\L}{\mathbb{L}}
\newcommand{\U}{\mathbb{U}}
\newcommand{\st}{^*}
\newcommand{\tq}{\middle|}
\newcommand{\inv}{^{-1}}
\newcommand{\E}{\mathbb{E}}
\newcommand{\F}{\mathbb{F}}
\newcommand{\V}{\mathbb{V}}
\renewcommand{\ker}{\text{Ker}}
\newcommand{\im}{\text{Im}}
\newcommand{\card}{\text{Card}}
\newcommand{\Tr}{\text{Tr}}

\title{TD Euclidiens\\ Exercice 33}
\date{}
\begin{document}
\maketitle

\section{Énoncé}
Soit $M\in\R_+$.
\begin{enumerate}
	\item Montrer que l'ensemble des matrices symétriques à spectre inclus dans $[0,M]$ est un compact de $\mathcal{S}_n(\R)$.
	\item $\star$ Montrer que $A\mapsto A^2$ est un homéomorphisme de l'ensemble des matrices symétriques positives dans lui-même.
\end{enumerate}
\section{Solution}
La deuxième question est difficile.
\begin{enumerate}
	\item Orthodiagonaliser, utiliser la compacité de $\mathcal{O}_n(\R)$ et de $[0,M]^n$ puis le fait que $\mathcal{S}_n(\R)$ est fermé.
	\item Par théorème de cours HP, cette application est bijective. Elle est de plus continue car polynomiale. On veut montrer que sa réciproque est continue. On prend $A_p$ qui tend vers $A$. On montre que toutes les valeurs d'adhérence de $(A_p^{1/2})$ sont nécessairement $A^{1/2}$ par unicité de la limite et compacité des ensembles mis en jeu (on a alors affaire à des fermés bornés car le spectre des matrices est borné) en utilisant la première question.
\end{enumerate}
\end{document}
